\section*{ID10008: Mathematical Framework: sin(φ) > 0 Rightarrow}
\paragraph{Metadata.}
PiID: 6619885782794 | RTEID: RTE:P38863.2866:T1.2029 | UUID: f6791492-0f1d-550a-abae-7e55e4e9db2c

\paragraph{Canonical Statement.}
This expression reveals the bifurcation mechanism. The sign of \$\textbackslash{}sin(\textbackslash{}phi)\$ changes at the equator: \textbackslash{}begin\{itemize\} \textbackslash{}item Northern Hemisphere: \$\textbackslash{}sin(\textbackslash{}phi) > 0 \textbackslash{}Rightarrow\$ rightward deflection \textbackslash{}item Equator: \$\textbackslash{}sin(\textbackslash{}phi) = 0 \textbackslash{}Rightarrow\$ no deflection \textbackslash{}item Southern Hemisphere: \$\textbackslash{}sin(\textbackslash{}phi) < 0 \textbackslash{}Rightarrow\$ leftward deflection \textbackslash{}end\{itemize\}

\paragraph{Canonical Equation.}
\[
\sin(\phi) > 0 \Rightarrow
\]

\paragraph{Dictionary.}
This expression reveals the bifurcation mechanism [ID10008]

\paragraph{Encyclopedia.}
Mathematical Framework: sin(φ) > 0 Rightarrow is an inequality / bound, written sin(φ) > 0 Rightarrow. It states an inequality or bound that constrains the admissible range of the quantity. It is built from φ. Within the Trawin Zero Point registry it serves as one element of the framework's mathematical narrative.
